\subsection{正常一次境界の非零 $y^2$ 項：稠密開集合の排除}
\label{subsec:cubic-third-root-unit-normal-boundary-open}
\begin{theorem}[正常一次境界商の開集合障害]
\label{thm:cubic-third-root-unit-normal-boundary-open}
前定理の六係数表示で $w\ne0$ とする。
生成元のscalar変更で $w=1$ とした後、以下の多項式を定義する：
\begin{align*}
 L&=D^3+\omega E^3+\omega^2,\\
 C_j&=D^2E+\omega^{-j}D+\omega^j E^2\quad(j=0,1,2),\\
 \Delta&=DE(D^3-E^3)(E^3-1)
 (D^3+\omega E^3)(D^3+\omega^2E^3)C_0C_1C_2L.
\end{align*}
$\Delta\ne0$ ならH3切断は存在しない。これは非空稠密開集合の排除であり、
$\Delta=0$ の曲線又は $w=0$ の全排除ではない。
\end{theorem}
\begin{proof}
$y\mapsto s y,b\mapsto r b,c\mapsto r c$ とし、$\gamma=r/s$ と置くと、
三生成元境界商の新係数は $D'=\gamma D,E'=\gamma E,w'=\gamma^3w$。
従って代数閉体上で $w'=1$ にでき、四生成元の六関係の形も保たれる。
$x$ は正常正則なので、$S=B/(x)$ は三関係
\[
 Dyb+Eby+c^2=0,\qquad Eyc+Dcy+b^2=0,\qquad
 Dbc+Ecb+y^2=0
\tag{*}
\]
を持つ。
H3の候補 $G$ の像 $\overline G$ が零なら $G\in xB_2$ で非整域切断となる。
従って $\overline G\ne0$ が必要であり、これは中山固定部分の中で
\[
 \overline Gy=y\overline G,\quad
 \overline Gb=\omega^2b\overline G,\quad
 \overline Gc=\omega c\overline G
\tag{**}
\]
を満たさなければならない。

元の $(*)$ の三次・四次関係挿入行列を
$\mathbb Q(\omega)(D,E)$ 上で消去する。
それぞれrank17,66であり、商次元は10,15。
全83個の関係pivotの分子・分母の非定数因子は $\Delta$ の因子に限る。
これは因数分解の推測でなく、各分子・分母を $\Delta$ とのgcdで逐次割り切り、
残りが定数となることを公開検証器で確かめる。
従って $\Delta\ne0$ の全点で同じ商基底を使える。
中山固定三次の自由基底は $b^3,bcy,cby,c^3$ の四つ。

$(**)$ の全係数行列を作り、最初の三つのpivotで三係数を消去する。
これらのpivotも分子・分母が $\Delta$ 上だけで消える。
残る第四係数に対して九つの残差 $r_i\in\mathbb Q(\omega)(D,E)$ が得られる。
$r_i$ の分母も全て $\Delta$ 上だけで消える。
その分子を $n_i$ とし、$I=(n_1,\ldots,n_9)$ と置く。
元の九分子から計算したGroebner基底でexactに
\[
 \bigl(E(E^3-1)L\bigr)^3\in I
\tag{***}
\]
となる。基底の全係数、元行列からの再生、$(***)$ の零剰余をcertificateと検証器に保持する。
$\Delta\ne0$ なら $(***)$ の左辺は非零なので、全 $r_i$ が同時に零になることはない。
従って第四係数も零、正常核は零。
これは一点のrank4を一般化した議論ではなく、全pivot単元性と残差イデアルによる開集合全体の証明である。
$u,v,Z$ は $(*)$ に現れないので、それらの全値を同時に排除する。
なお $D=2,E=3$ では $\Delta\ne0$ であり、開集合は非空。
\end{proof}
